\centerline{\bf \Huge Всероссная тренировка I}
\centerline{\bf \Huge }

\begin{flushright}
        \scriptsize{вот бы так же на всероссе было}
\end{flushright}

\problem{1}
В английском королевстве есть 30 достойных рыцарей. Каждый из них проводит за год 82 рыцарских поединка с другими рыцарями в регулярном рыцарском турнире. Можно ли их разделить на две фракции Ланкастеров и Йорков и составить расписание поединков так, чтобы бои между рыцарями из разных фракций составляли ровно половину от общего числа поединков?


\problem{2}
Окружность $\omega$ касается сторон угла $B A C$ в точках $B$ и $C$. Прямая $\ell$ пересекает отрезки $A B$ и $A C$ в точках $K$ и $L$ соответственно. Окружность $\omega$ пересекает $\ell$ в точках $P$ и $Q$. Точки $S$ и $T$ выбраны на отрезке $B C$ так, что $K S \| A C$ и $L T \| A B$. Докажите, что точки $P, Q, S$ и $T$ лежат на одной окружности.


\problem{3}
Диагонали $A C$ и $B D$ вписанного четырехугольника $A B C D$ пересекаются в точке $P$. Точка $Q$ выбрана на отрезке $B C$ так, что $P Q \perp A C$. Докажите, что прямая, проходящая через центры окружностей, описанных около треугольников $A P D$ и $B Q D$, параллельна прямой $A D$.

\problem{4}
В треугольнике $ABC$ проведены чевианы $AA_1$, $BB_1$ и $CC_1$, пересекающиеся в одной точке. Прямая $B_1C_1$ пересекает окружность $(ABC)$ в точках $P$ и $Q$. Докажите, что окружность $(A_1PQ)$ проходит через середину отрезка $BC$.